\chapter*{List of aids used}
\addcontentsline{toc}{chapter}{List of aids used}

\begingroup
\RedeclareSectionCommand[beforeskip=1.2\baselineskip,afterskip=0.4\baselineskip]{section}

\section*{Literature access and reference verification}

arXiv, ACL Anthology, OpenReview and publisher websites were used to access
research papers and check publication details. Individual papers are cited
where their content is used and listed in the bibliography.

\section*{Published method and code}

The method reimplemented here is that of Farquhar et
al.~\cite{semantic_entropy_nature_2024}. The reference implementation is the
authors' released code~\cite{kossen_semantic_uncertainty_code_2024}, which was
used as the specification for clustering, prompting, P(True) scoring and
question selection. Deviations from it are recorded in
Section~\ref{sec:deviation-register}.
The reference implementation was accessed through GitHub; its archived
release on Zenodo is identified by the accompanying code citation.

\section*{Datasets}

TriviaQA~\cite{joshi_triviaqa_2017}, NQ-Open~\cite{lee_nq_open_2019} and
SVAMP~\cite{patel_svamp_2021}, prepared using a reimplementation of the
released loading and selection procedure. The biography prompts follow the
FActScore subject format~\cite{min_factscore_2023}.
Hugging Face Hub was used to obtain model checkpoints and datasets and to
consult model cards.

\section*{Models}

\begin{itemize}
  \raggedright
  \item \emph{Answer generation}, at the pinned revisions given in
    Table~\ref{tab:generators}: \texttt{Mistral-7B-Instruct-v0.3},
    \texttt{Llama-3.1-8B-Instruct} and \texttt{Llama-3.1-70B-Instruct}.
  \item \emph{Entailment}: \texttt{microsoft/deberta-v2-xlarge-mnli},
    \texttt{cross-encoder/nli-deberta-v3-large},
    \texttt{cross-encoder/nli-deberta-v3-base} and
    \texttt{Qwen/Qwen2.5-72B-Instruct}.
  \item \emph{Correctness judging}: \texttt{Qwen/Qwen2.5-72B-Instruct} and
    \texttt{Llama-3.1-70B-Instruct}.
\end{itemize}

\section*{Software}

Python with \texttt{PyTorch}, \texttt{Transformers}, \texttt{vLLM},
\texttt{NumPy}, \texttt{pandas}, \texttt{scikit-learn}, \texttt{Matplotlib},
\texttt{datasets} and \texttt{pytest}. Environments were managed with
\texttt{uv}, revision control with \texttt{Git}. The thesis was typeset in
\LaTeX\ using the \texttt{kaobook} class, with \texttt{Biber} and
\texttt{biblatex} for the bibliography. Builds were managed with
\texttt{latexmk}.

\section*{Computing resources}

The NHR@FAU Alex GPU cluster for answer generation and language-model judging,
and the NHR@FAU Fritz cluster for post-hoc cross-encoder clustering, scheduled
with \texttt{Slurm}.

\section*{Artificial-intelligence assistants}

Claude Code (Anthropic) and Codex (OpenAI) were both used as coding assistants
in the accompanying repository and to generate the figures and tables in this
thesis from the recorded measurements. As stated in the declaration, these
tools were not used
as a source of claims, results or citations. Every reported number was checked
against the recorded measurements.
\par
\endgroup
